\documentclass[conference]{IEEEtran}
\usepackage[utf8]{inputenc}
\usepackage[T1]{fontenc}
\usepackage{cite}
\usepackage{amsmath,amssymb}
\usepackage{graphicx}
\usepackage{booktabs}
\usepackage{url}

\title{Do Generate--Verify--Repair Guards Improve LLM NetOps Outputs? A Paired Emulator Study}

\author{
\IEEEauthorblockN{Autonomous AI Researcher}
\IEEEauthorblockA{CNSM 2026 Agentic AI Researcher Track}
}

\begin{document}

\maketitle

\begin{abstract}
We preregistered and executed a paired confirmatory experiment (n=200 natural-language NetOps intents; study\_id cnd-001-5f3a) that compares ungated LLM generation to a bounded generate--verify--repair (GVR) pipeline applied to a single shared LLM candidate per intent. Generation used gpt-5-mini and the guarded arm applied a deterministic three-stage validator and a deterministic, rule-based repair operator with repair\_budget <=1 (no extra LLM calls). Artifact-backed primary outcomes: baseline successes = 199/200 (0.995), guarded = 200/200 (1.000); paired contingency n11=199, n10=1, n01=0, n00=0; paired difference (guarded - baseline) = 0.005. The preregistered exact two-sided McNemar (exact binomial on discordant pairs) yields p = 1.0; paired-bootstrap (10,000 resamples, seed 1729) 95\% percentile CI = [0.0, 0.015]. The preregistered 0.15 absolute-effect target is excluded by the observed CI because the realized baseline was near-ceiling. We provide concise, auditable descriptions of the validator and repair operator, execution accounting (model\_calls\_used = 201; deterministic repair invocations = 1), exact-test justification appropriate for small discordant counts, and operationally grounded limitations. All prompts, provider traces, validator traces, repair traces (the single invoked repair trace), and analysis artifacts are archived in the study evidence bundle.
\end{abstract}

\section{Introduction}
Automating routine network-operations (NetOps) changes with large language models (LLMs) promises reduced manual effort but raises an operational-safety question: do LLM-produced configuration artifacts execute correctly when applied to control planes? Prior work demonstrates intent-to-configuration prototypes and documents structured-output failures (missing commands, misordering, protected-field modification) and proposes mitigations such as retrieval-augmentation, constrained decoding, verifier-assisted repair, and multi-stage tool-augmentation. Despite these proposals, there is limited execution-grounded evidence that quantifies the marginal benefit of applying a bounded deterministic repair to the same initial LLM candidate (a low-hosted-call operational estimand).

We preregistered and executed a paired experiment (study\_id cnd-001-5f3a) that isolates the ''repair-on-same-candidate'' question by sharing a single initial LLM candidate across both arms (baseline ungated acceptance and guarded generate--verify--repair). The shared-candidate estimand is operationally relevant when hosted-model call budgets are constrained. Our contributions are: (1) a preregistered, artifact-backed paired evaluation under the shared-candidate estimand; (2) concise algorithmic descriptions of the deterministic three-stage validator and the deterministic, rule-based repair operator used in the guarded pipeline; and (3) complete execution accounting and exact-test justification for small-discordant-pair inference.

\section{Related Work}
Research on LLMs for NetOps has rapidly diversified along three complementary directions: intent-to-configuration translation and co‑pilot prototypes, structured-output mitigation techniques, and formal/deterministic validators that can be integrated into tool-augmented pipelines. Prior prototype and survey work document practical promise---natural-language intents can be translated into device-level commands and assist operators in configuration generation---but also recurrent failure modes such as missing required commands, incorrect ordering, and unintended modification of protected fields (intent‑translation prototypes summarized in recent literature) [doi:10.1002/nem.2313; https://openalex.org/W7161165783].

Mitigation efforts borrow from nearby domains. Retrieval-augmented generation (RAG), live schema grounding, and constrained decoding have shown material reductions in hallucinated or structurally invalid outputs in NL2SQL and document-grounded tasks; these approaches are natural candidates for NetOps but require careful adaptation to device CLI dialects, workflow policies, and operational constraints [doi:10.2139/ssrn.6909831]. Similarly, constrained or grammar-aware decoding methods and deterministic post-hoc canonicalizers reduce syntactic errors but do not by themselves guarantee correct control‑plane behavior when applied to live or emulated networks.

Deterministic verification and formal-methods work (NetKAT variants, weighted NetKAT and ensemble verifiers) provide a complementary assurance axis by checking reachability, isolation, and quantitative properties of candidate configurations; these techniques can produce machine-readable violation codes that are actionable for automated repair or operator review [doi:10.1145/3808318]. However, verification tools face two practical constraints when applied in NetOps pipelines: (1) fidelity to dynamic control-plane phenomena (timing, convergence, vendor-specific idiosyncrasies) and (2) scalability when reasoning over larger multi-device changes. These limits motivate hybrid architectures that combine fast deterministic checks with auditable traces and bounded repair operations.

Security- and safety-oriented evaluations emphasize adversarial risks: human‑readable jailbreak prompts and adversarial-context attacks remain an effective threat class for unconstrained generative pipelines and therefore constitute an operational risk for NetOps automation unless mitigations and guards are deployed [doi:10.1145/3815174]. Architectural discussions of agentic NetOps emphasize that production reliability depends on the surrounding assurance machinery---auditable evidence traces, bounded tool use, rollback policies, and independent validators---rather than on raw model capability alone (see recent agentic NetOps position work) [https://openalex.org/W7161165783].

Across these literatures the recurring gap is empirical and execution-grounded: many studies show reduced failure modes in isolated subcomponents (generation, decoding, or verification) but there are relatively few reproducible, artifact-backed measurements of end-to-end generate--validate--repair pipelines under practical constraints such as limited hosted-model calls, deterministic repair budgets, and auditable validator traces. Our study targets this specific gap by measuring the marginal effect of a bounded deterministic repair applied to the same initial LLM candidate under an explicit shared-generation budget; this positions our work as an execution-grounded complement to prior prototype, mitigation, and verification literature [doi:10.1002/nem.2313; doi:10.2139/ssrn.6909831; doi:10.1145/3808318; https://openalex.org/W7161165783; doi:10.1145/3815174].

\section{Methodology}
Preregistration and high-level design: We preregistered a paired, shared-initial-generation confirmatory protocol (study\_id cnd-001-5f3a). For each of n = 200 curated natural-language intents the hosted model produced one initial candidate configuration which was then scored directly as the baseline arm and also passed to the guarded arm. The guarded arm applied a deterministic validator; when the validator reported a repairable violation the pipeline invoked a deterministic, rule-based repair operator (repair\_budget <= 1 for the confirmatory run). The primary estimand is the paired success-rate difference (guarded - baseline); the preregistered primary test is the exact two-sided McNemar implementation (exact-binomial on discordant pairs). Planned maximum hosted-model calls = 400; realized model\_calls\_used = 201 (shared\_initial\_generation = 200; deterministic repair-stage calls = 1). The execution environment combined Dockerized Mininet+FRR with deterministic\_netops\_validator\_v5 as the canonical validator.

Task bank and operational scoring: The confirmatory bank contains 200 intent--expected-configuration pairs created by a deterministic task generator. Each task payload includes initial\_state, expected\_final\_state, required\_sequence (ordered per-task commands the authoritatively required), workflow policies (e.g., make-before-break, do-not-touch protected interfaces), and difficulty metadata. The scoring rule (full\_intent\_constraint\_validation) requires that generated configurations (a) parse syntactically into canonical CLI-like statements, (b) include the task's required\_sequence (order and presence), (c) respect dependency and group-activity policies (e.g., minimum-active constraints), (d) preserve protected-interface values, and (e) produce an emulator-observed final\_state equal to expected\_final\_state when applied in the deterministic emulator. Provider/API generation failures are treated per the preregistered missingness plan as arm failures; none occurred in the confirmatory run.

Deterministic validator: The validator deterministically executes three ordered stages and emits a machine-readable trace that the scoring and repair logic consume. Stage 1 -- syntactic parsing and canonical normalization: tokenizes CLI-like lines, canonicalizes variant synonyms, and produces normalized\_configuration; syntax failures yield a SYNTAX\_PARSE\_ERROR and block further deterministic repair. Stage 2 -- structural and workflow checks: the validator deterministically verifies presence and ordering of required\_sequence entries and enforces dependency rules and group policies; it emits deterministic violation codes such as MISSING\_REQUIRED\_COMMAND, DEPENDENCY\_VIOLATION, PROTECTED\_INTERFACE\_MODIFICATION, and TRANSIENT\_ORDER\_VIOLATION. Stage 3 -- deterministic emulator application: the normalized\_configuration is applied to the emulator to produce observed final\_state, and the validator reports validation\_after.valid based on exact equality with expected\_final\_state. The validator trace fields archived per episode (and used by analysis) include parsed\_command\_count, observed\_touched\_field\_count, normalized\_configuration, violation\_count, violation\_codes, validator\_id and validator\_version, and validation\_before/after summaries. These artifacts support auditable, post-hoc inspection of every failure mode.

Repair operator: The preregistered repair operator is a deterministic, rule-based transformer invoked only when the validator reports violation\_count > 0 and the violation family is in a preregistered repairable set. The operator enforces the post-lock constraints: (i) no additional LLM calls (repairs are implemented by deterministic code); (ii) bounded budget (<=1 repair attempt per episode in confirmatory runs); and (iii) unambiguous mapping requirement (repairs are applied only when the mapping from violation code to transformation is unique given the task payload). Repair decisions are computed from the task manifest (not inferred by an LLM): the operator consults required\_sequence, initial\_state, expected\_final\_state, and workflow\_policies to deterministically generate corrections. The operator implements three canonical deterministic transformations:
- Insertion: when a MISSING\_REQUIRED\_COMMAND violation exists and the missing step appears in the task's required\_sequence, the operator inserts the missing command(s) at canonical positions defined by the required\_sequence. Positioning rules are deterministic: insert before the first command that references the same interface/field group or at the canonical step index derived from the required\_sequence ordering in the manifest.
- Reordering: for DEPENDENCY\_VIOLATION or TRANSIENT\_ORDER\_VIOLATION the operator performs a stable reordering of contiguous, annotated command blocks to satisfy dependency constraints while preserving intra-block lexical order. The reordering algorithm deterministically identifies minimal block moves that resolve the dependency according to the manifest's dependency rules.
- Preservation enforcement (redaction/restore): for PROTECTED\_INTERFACE\_MODIFICATION or when preserve\_unspecified\_state would be violated, the operator deterministically replaces the offending command(s) with the original value obtained from initial\_state (or, if required by workflow\_policies, with the canonical preserved value from expected\_final\_state).

Ambiguity and nonrepairable cases: If a violation admits multiple equally plausible deterministic repairs (ambiguous insertion indices, multiple candidate values), the operator abstains (no change) and the episode remains a guarded failure; ambiguous families and policy-contradictory violations are flagged in the repair\_trace and counted in analysis as nonrepairable. All repair decisions and the resulting normalized\_configuration are recorded in a repair trace archived with the run artifacts. In the confirmatory run the deterministic repair operator was invoked exactly once and that single repair converted its episode to a validator-passing configuration.

Statistical analysis: The primary confirmatory test is McNemar's paired binary test implemented in its exact-binomial (exact McNemar) form appropriate for small discordant counts. Let n10 be the number of pairs where guarded = 1 and baseline = 0, and n01 where baseline = 1 and guarded = 0; n = n10 + n01 is the total number of discordant pairs. The exact two-sided p-value is computed from the binomial distribution on n trials assuming p = 0.5 (symmetry under the null): the two-sided p equals twice the smaller of the lower and upper tail probabilities for observing k = n10 successes under Binomial(n, 0.5), truncated at 1.0. For our observed counts (n10 = 1, n01 = 0, n = 1) this exact-binomial two-sided computation yields p = 1.0. Complementary uncertainty quantification uses a paired-bootstrap percentile 95\% CI for the guarded - baseline proportion difference with 10,000 resamples (bootstrap\_seed = 1729) as preregistered; these parameters and the bootstrap result are archived in analysis artifacts. The preregistered power analysis assumed a baseline \ensuremath{\approx}0.5 and targeted a 15 percentage-point absolute effect at n = 200; the realized near-ceiling baseline substantially reduces the expected discordant counts and consequent sensitivity for that target effect size. All analysis scripts, seeds, resample counts, and the paired contingency CSV are archived in the evidence bundle for reproducibility.

Operational accounting and failure rules: Provider/API transient errors were treated per the preregistered missingness plan (single automatic retry for transient timeouts with 60s backoff); no such retries were consumed in the confirmatory run. Episodes failing syntactic parse are unrepaired by design (repair operator requires a valid normalized\_configuration); such parse failures were zero in the confirmatory bank. Episodes where repair abstains are reported as guarded failures and counted in the paired contingency table. The deterministic pipeline's audit traces (per-episode validator traces and any repair traces) preserve the exact sequence of decisions for post-hoc inspection and are referenced by analysis outputs.

\section{Results}
Execution and primary outcomes: The confirmatory run completed all planned episodes (400 episodes = 200 pairs) with no provider/API failures. Artifact-backed primary counts: baseline\_success\_count = 199/200 (0.995) and guarded\_success\_count = 200/200 (1.0). The paired contingency table is n11 = 199, n10 = 1 (guarded-only), n01 = 0 (baseline-only), n00 = 0; paired difference (guarded - baseline) = 0.005. The preregistered exact two-sided McNemar (exact binomial on discordant pairs) yields p = 1.0 for the observed ordering (n = 1). The paired-bootstrap (10,000 resamples, seed 1729) 95\% percentile CI is [0.0, 0.015]. All per-episode JSONL scoring artifacts and the analysis result JSON are archived.

Repair and failure accounting: validation\_before.valid == false occurred in exactly one confirmatory episode; deterministic repair invocations = 1 and that single repair converted the failing candidate into a validator-passing final configuration, producing the sole guarded-only success (n10 = 1). Syntactic/parse failures across the confirmatory set = 0. Provider-call audit: hosted\_model\_call\_event\_count = 201, completed\_hosted\_model\_call\_event\_count = 201, cache\_miss\_count = 201, stage\_counts show shared\_initial\_generation = 200 and repair = 1. No episodes were excluded for contamination; missingness summary reports complete\_pairs = 200 and zero unscorable episodes.

Representative exemplars and difficulty context: we archive six representative exemplars with per-example difficulty labels (medium = 1, hard = 2, very\_hard = 3) to illustrate typical intent patterns including shutdown\_configure\_restore, make-before-break uplink switchover, and paired atomic MTU changes. The full per-episode difficulty distribution for all 200 confirmatory intents is recorded in the archived task manifest (execution/task\_manifest.jsonl) and is available for re-analysis; because aggregated counts are stored in that manifest we reference it for exhaustive difficulty counts rather than enumerating the full distribution inline.

Sensitivity and interpretive statistics: The preregistered 0.15 absolute-effect target was chosen under a baseline \ensuremath{\approx}0.5 assumption and is excluded by the observed bootstrap CI; the near-ceiling baseline compresses attainable discordant counts and reduces power to detect large absolute improvements under the shared-candidate estimand. Practically, deterministic repair-on-same-candidate produced a single marginal improvement in this curated confirmatory bank.

\section{Discussion}
We expand the operational interpretation and artifact-grounded diagnostics to help practitioners evaluate when a guarded, same-candidate GVR gate is worth deploying.

1) Failure-mode diagnostics and repair yield. The archived validator traces show that across the 200 confirmatory intents the validator emitted zero SYNTAX\_PARSE\_ERROR events and only one episode with violation\_count>0 that the deterministic repair operator classified as repairable under the preregistered mapping. That single episode illustrates two points: (a) the dominant failure modes in this curated bank were not low-level parsing errors but rather missing or misordered commands that the validator could deterministically resolve when the task payload supplied an explicit required\_sequence; (b) the deterministic repair yield is therefore highly dependent on task design---when intents include explicit required\_sequence elements the operator can deterministically insert or reorder commands with high confidence. The repair trace for the single invoked repair records the deterministic transformation (insertion of a missing shutdown/restore pair at canonical positions) and the validator re-run that produced validation\_after.valid = true; both the trace and provider-call audit are archived for inspection.

2) Cost--yield trade-offs under hosted-call budgets. Execution accounting (model\_calls\_used = 201; shared\_initial\_generation = 200; repair-stage calls = 1) quantifies the low hosted-call cost of the shared-candidate estimand: a single extra deterministic repair invocation (no additional LLM calls in our run) sufficed to reach 100\% guarded success in this bank. In contexts where provider cost or throttling is binding, this design gives nearly full assurance at minimal extra call cost. By contrast, an independent-resampling estimand (two or more independent LLM generations per intent) would increase hosted calls proportionally; the archived per-episode shared\_initial response SHA values facilitate counterfactual replays to estimate expected incremental yield from resampling without incurring provider costs by simulating alternative-generation distributions from stored LLM completions.

3) Interpreting the near-ceiling baseline and statistical sensitivity. The near-ceiling baseline (199/200) produced only one discordant pair (n = n10+n01 = 1), which forces the exact-binomial McNemar test to reflect extreme small-n behavior (two-sided p = 1.0 under the observed ordering). The exact McNemar formulation used (binomial test on discordant pairs) is appropriate for such small n because the common large-sample chi-squared approximation is invalid. The paired-bootstrap CI (10,000 resamples; seed = 1729) provides a complementary finite-sample uncertainty interval [0.0,0.015] for the absolute paired difference. Practitioners should therefore interpret our null inferential finding not as evidence that repairs never help, but as evidence that for this curated task bank and single-model generation the marginal improvement achievable by a conservative deterministic repair-on-the-same-candidate is small and auditable. Where baseline error rates are higher, the same guarded machinery would produce larger discordant counts and greater statistical power to detect meaningful improvements.

4) Validator design implications. The three-stage deterministic validator demonstrates operational value beyond pass/fail classification: it emits structured violation codes (MISSING\_REQUIRED\_COMMAND, DEPENDENCY\_VIOLATION, PROTECTED\_INTERFACE\_MODIFICATION) and quantitative difficulty metadata (parsed\_command\_count, observed\_touched\_field\_count, required\_command\_count). These structured outputs enabled our repair operator to apply targeted deterministic actions only when the mapping from violation code to transformation was unambiguous. For production pipelines, we recommend exposing the validator's violation taxonomy to policy layers (e.g., escalation thresholds for HUMAN\_REVIEW when violations fall outside the repairable set) and recording violation-code histograms over time for operational monitoring; the archived analysis/contamination\_summary.csv and per-episode validator\_traces provide the raw counts needed to implement such monitoring.

5) Emulation fidelity and external generalizability. The deterministic emulator applied the normalized\_configuration to a Dockerized Mininet+FRR environment and compared observed final\_state to expected\_final\_state. While this provides deterministic, reproducible validation for the control-plane state modeled, the emulator does not model device-specific firmware behaviors, asynchronous control-plane convergence timing (e.g., BGP timers), or vendor command dialect subtleties. The single repair that succeeded did so by addressing a syntactic/ordering gap that is emulator-visible; failures attributable to runtime dynamics would not necessarily be repairable by deterministic static transforms. Users should treat emulator success as necessary but not sufficient evidence for safe production deployment without device-in-the-loop testing.

6) When a guarded same-candidate gate is most appropriate. Based on artifact-grounded evidence, we recommend the shared-candidate GVR gate when: (a) hosted-call budgets are constrained; (b) intents are well-structured with explicit required\_sequence or workflow policies; and (c) the validator's repairable-family mapping covers the common failure modes (insertions/reorderings/preservation). When intents are open-ended, diverse in CLI dialect, or when high baseline error rates are observed in validator logs, independent resampling or LLM-invoked repair loops---despite higher provider cost---may be necessary. The execution manifest and provider\_call\_audit in the archived bundle offer concrete accounting numbers for such cost--benefit calculations.

7) Artifact-grounded reproducibility and next-step diagnostics. We provide all artifacts required to reproduce the run and to run targeted diagnostic re-analyses: per-episode shared\_initial responses (execution/responses/*), validator traces (execution/scoring/*.json), repair trace (single invoked repair), and the full task manifest (execution/task\_manifest.jsonl; SHA-256: eb2f045f\allowbreak{}8a6214eb\allowbreak{}6d814e16\allowbreak{}f5d95e01\allowbreak{}20bd8997\allowbreak{}3951b6fe\allowbreak{}14084bda\allowbreak{}4e3efb4a). Researchers can re-run paired analyses with alternate estimands (independent-resampling) or simulate resampling yield by sampling from archived candidate pools. The archived analysis/paired\_contingency\_table.csv and analysis/paired\_difference\_figure.svg provide the exact contingency counts and paired-difference diagnostics used in this paper.

Taken together, these diagnostics show that deterministic validators paired with conservative, auditable deterministic repairs can provide a low-cost assurance gate that is effective when baseline generator quality is high and intents are structured. The approach is not a universal substitute for higher-cost remediation strategies (resampling, LLM-guided repair, or human-in-the-loop review) in settings with higher baseline error rates, multivendor CLI diversity, or dynamic runtime failure modes.

\section{Limitations}
\begin{itemize}
\item Synthetic, curated confirmatory task bank focused on small single- and multi-device changes; not comprehensive of production NetOps workloads (multivendor CLI dialects, hardware idiosyncrasies, dynamic control-plane timing).
\item Single LLM (gpt-5-mini) evaluated; findings do not imply multi-model generality.
\item Deterministic emulator + validator model control-plane behavior but cannot fully capture real-world runtime dynamics such as BGP convergence timing or vendor-specific runtime effects.
\item Preregistered repair budget limited to a single deterministic repair iteration; different repair strategies or additional iterations might yield different outcomes.
\item Shared-initial-generation design reduces model-call cost and isolates repair-on-same-candidate effectiveness but reduces external generalizability compared with independent-resampling estimands.
\item Near-ceiling baseline success limited statistical power to analyze failure-mode distributions in the confirmatory set; exploratory failure-taxonomy conclusions are therefore constrained.
\item Adversarial or out-of-distribution intents (e.g., jailbreaks, malformed ambiguous intents) were not included in the confirmatory set and require separate evaluation.
\item Full per-task difficulty label counts for all 200 confirmatory intents are recorded in execution/task\_manifest.jsonl in the archived evidence bundle; the manuscript references that manifest rather than enumerating the full 200-line distribution inline to preserve concise presentation while preserving auditable reproducibility.
\end{itemize}

\section{Conclusion}
We executed an autonomy-locked, preregistered paired experiment to measure whether a bounded generate--verify--repair guard improves emulator-executable success for LLM-generated NetOps configurations under a shared-candidate generation budget. Ungated gpt-5-mini generation succeeded in 199/200 confirmatory cases; the deterministic guarded pipeline succeeded in 200/200. The observed paired difference = 0.005 (exact McNemar two-sided p = 1.0; paired-bootstrap 95\% CI [0.0, 0.015]) does not support the preregistered 0.15 absolute-effect target because the baseline was near-ceiling. For this estimand and task distribution, deterministic repair-on-same-candidate yielded negligible marginal improvement, but deterministic validators remain valuable, auditable assurance gates for operational NetOps pipelines. Full artifacts and analysis code are archived with the study for independent audit and re-analysis.

\section*{Disclosure Statement}
This study was executed autonomously after an autonomy lock defined by the initial master prompt archived at provenance/master\_prompt.txt (SHA-256: 1872df1e\allowbreak{}1805d2d9\allowbreak{}6940456c\allowbreak{}a016bd66\allowbreak{}5d1d5196\allowbreak{}add77f5a\allowbreak{}cdf1582b\allowbreak{}b39b15ba). Human involvement prior to the lock was limited to selecting the topic family (Generative AI and LLMs for NetOps) and approving the master prompt. No manual human editing of technical content, figures, captions, references, results, or manuscript text occurred after the lock. The autonomous pipeline performed literature search and verification, research-gap selection, preregistration (study\_id cnd-001-5f3a), code generation, experiment execution, artifact capture, analysis, internal critique, peer-review checks, and manuscript drafting.

Models and tooling: generation requested gpt-5-mini (declared\_model\_version in execution/model\_configuration.json). Validation used deterministic\_netops\_validator\_v5 implementing syntactic parsing, structural/workflow checks, and deterministic emulator application. Repairs in the confirmatory run were applied by deterministic code only (no additional LLM calls). The execution environment used Dockerized Mininet+FRR images orchestrated by adapter family hosted\_netops\_gvr\_v1. Preregistered and enforced settings (not altered post-lock) include: primary estimand = paired\_success\_rate\_difference\_guarded\_minus\_baseline; confirmatory sample = n=200 paired intents; maximum model-call budget = 400; repair budget <=1 deterministic repair iteration; primary analysis = exact McNemar two-sided (exact-binomial on discordant pairs); bootstrap CIs = 10,000 resamples, seed = 1729. Execution accounting (artifact-backed): the run completed 400 episodes with model\_calls\_used = 201 (shared\_initial\_generation = 200; repair-stage calls = 1). All prompts, provider-call traces, responses, validator traces, repair traces (the single invoked repair trace), analysis artifacts, and the execution manifest are archived in the study evidence bundle.

\begin{thebibliography}{99}
\bibitem{ref1} Nguyen Tu, Sukhyun Nam, James Won‐Ki Hong, "Intent‐Based Network Configuration Using Large Language Models", 2024, 10.1002/nem.2313
\bibitem{ref2} Emmanuel Suárez Acevedo, Tiago Ferreira, Kevin Batz, Oliver Bøving, Nate Foster, Alexandra Silva, "Weighted NetKAT: A Programming Language for Quantitative Network Verification", 2026, 10.1145/3808318
\bibitem{ref3} Sanjay Mishra, Divya Chukkapalli, Ganesh R. Naik, "Schema-Aware Localisation (SAL): Live Schema Grounding and Hallucination Validation for Oracle NL2SQL", 2026, 10.2139/ssrn.6909831
\bibitem{ref4} Muhammad Bilal, Jon Crowcroft, Ruizhi Wang, Xiaolong Xu, Schahram Dustdar, "Large Language Models for Agentic NetOps and AIOps: Architectures, Evaluation, and Safety", 2026, 10.48550/arxiv.2605.12729
\bibitem{ref5} Nilanjana Das, Edward Raff, Aman Chadha, Manas Gaur, "Human-Readable Adversarial Prompts: An Investigation into LLM Vulnerabilities Using Situational Context", 2026, 10.1145/3815174
\end{thebibliography}

\end{document}
